\documentclass[12pt,a4paper]{article}

% User-provided visual/layout reference imported for v8.0.0 Template Authority.
% This file is a reference exemplar only; executable competition output continues to use cumcmthesis.

\usepackage{geometry}
\geometry{left=2.5cm,right=2.5cm,top=2.5cm,bottom=2.5cm,a4paper}
\usepackage[UTF8]{ctex}
\usepackage{amsmath,amssymb}
\usepackage{graphicx}
\usepackage{float}
\usepackage{fancyhdr}
\usepackage{titlesec}
\usepackage{parskip}
\usepackage{caption}
\usepackage{subcaption}
\usepackage{booktabs}
\usepackage{array}
\usepackage{enumitem}
\usepackage{hyperref}
\usepackage{xcolor}
\usepackage{appendix}
\usepackage{listings}

\lstset{
  language=Python,
  basicstyle=\ttfamily\small,
  keywordstyle=\color{blue}\bfseries,
  stringstyle=\color{red},
  commentstyle=\color{green!40!black},
  numbers=left,
  numberstyle=\tiny\color{gray},
  stepnumber=1,
  numbersep=5pt,
  showspaces=false,
  showstringspaces=false,
  showtabs=false,
  frame=single,
  tabsize=4,
  captionpos=b,
  breaklines=true,
  breakatwhitespace=true
}

\pagestyle{fancy}
\fancyhf{}
\fancyfoot[C]{\thepage}
\renewcommand{\headrulewidth}{0pt}
\renewcommand{\footrulewidth}{0pt}

\setlength{\parindent}{2em}
\setlength{\parskip}{0pt}
\renewcommand{\baselinestretch}{1.25}

\titleformat{\section}{\centering\bfseries\Large}{\chinese{section}、}{0em}{}
\titleformat{\subsection}{\bfseries\large}{\arabic{section}.\arabic{subsection}}{1em}{}
\titleformat{\subsubsection}{\bfseries\normalsize}{\arabic{section}.\arabic{subsection}.\arabic{subsubsection}}{1em}{}
\titlespacing{\section}{0pt}{0pt}{20pt}
\titlespacing{\subsection}{0pt}{10pt}{10pt}
\titlespacing{\subsubsection}{0pt}{8pt}{8pt}

\captionsetup{labelsep=space,font=normalsize,labelfont=bf}
\hypersetup{colorlinks=true,linkcolor=black,citecolor=black,urlcolor=black}

\title{\textbf{数学建模论文版式参考模板}}
\author{}
\date{}

\begin{document}
\maketitle
\thispagestyle{empty}

\begin{center}
{\Large\textbf{摘要}}
\end{center}

% 摘要正文占位。

\textbf{关键词：}关键词一；关键词二；关键词三

\newpage
\setcounter{page}{1}
\pagestyle{fancy}

\section{问题重述}
\subsection{问题背景}
\subsection{问题提出}

\section{问题分析}
\subsection{问题一的分析}
\subsection{问题二的分析}

\section{模型假设}
\begin{enumerate}[leftmargin=2em]
  \item 假设一；
  \item 假设二；
  \item 假设三。
\end{enumerate}

\section{符号说明}
\begin{table}[H]
\centering
\caption{主要符号说明}
\begin{tabular}{ccl}
\toprule
\textbf{符号} & \textbf{含义} & \textbf{单位/说明}\\
\midrule
$x$ & 示例变量 & --\\
\bottomrule
\end{tabular}
\end{table}

% 原始参考文件在此之后采用“模型的建立与求解”统一一级章。
% v8.0.0 CUMCM 正式模板不沿用该章节组织，而改为每个问题独立一级章：
% “问题X模型建立及求解”。这一文件仅保留版式、编号、摘要、表格、代码和附录等视觉参考。

\section{模型的建立与求解}
\subsection{问题一的模型}
\subsubsection{模型建立}
\subsubsection{模型求解}
\subsection{问题二的模型}
\subsubsection{模型建立}
\subsubsection{模型求解}

\section{结果分析与验证}

\section{模型的评价与改进}
\subsection{模型的优点}
\subsection{模型的不足与改进方向}

\section*{参考文献}

\newpage
\appendix
\titleformat{\section}{\centering\bfseries\Large}{附录\Alph{section}}{1em}{}
\titleformat{\subsection}{\bfseries\large}{\Alph{section}.\arabic{subsection}}{1em}{}
\section{支撑材料列表}
\section{程序代码}

\end{document}
